\documentclass[11pt]{article}
\usepackage[margin=25mm]{geometry}
\usepackage{fontspec,unicode-math}
\setmainfont{Latin Modern Roman}
\setmathfont{Latin Modern Math}
\usepackage{amsmath,amssymb,mathtools}
\usepackage{xurl,hyperref}
\hypersetup{colorlinks=true,urlcolor=blue,linkcolor=blue}
\setcounter{secnumdepth}{0}
\setlength{\parindent}{0pt}
\setlength{\parskip}{5pt}
\setlength{\emergencystretch}{3em}
\newcommand{\sref}[2]{\hyperlink{source:#1:#2}{[S#2]}}
\newcommand{\cataloguescope}[1]{\paragraph{Source question.}#1}
\newcommand{\floor}[1]{\left\lfloor#1\right\rfloor}
\begin{document}
\section{Global Margin Maximization beyond the Neural Tangent Regime}\label{statement}

\subsection{Definitions and mathematical statement}

\paragraph{Shared notation.}Standard number systems are denoted by $\mathbb N,\mathbb Z,\mathbb Q,\mathbb R,\mathbb C$; index ranges, computational models and specialized notation follow the cited problem source and the definitions below.


For binary labeled data $(x_i,y_i)$ with $y_i\in\{-1,1\}$, the margin of a score function $f_\theta$ on sample $i$ is $y_i f_\theta(x_i)$. A parameter-homogeneous network satisfies $f_{c\theta}(x)=c^d f_\theta(x)$ for $c>0$ and a fixed degree $d>0$. A maximum-margin objective maximizes the smallest training margin subject to the source parameter normalization, or equivalently minimizes its normalized parameter cost under unit-margin constraints when that equivalence applies. Gradient descent updates $\theta_{t+1}=\theta_t-\eta_t\nabla L(\theta_t)$. The neural-tangent regime approximates the score by its first-order parameter expansion at initialization.
\paragraph{Statement recorded in the supplied source.}
Under what conditions does gradient descent globally optimize a maximum-margin objective for deep homogeneous networks, and when does it move beyond the neural-tangent regime to learn stronger representations? \sref{statement}{2}
\paragraph{Scope qualification.}
Architecture, loss, step sizes, separability and the margin normalization must be stated in any proposed answer; homogeneity alone is not claimed to ensure a global optimum.
\paragraph{Source status.}
The supplied catalogue labels this record open. Its source-status metadata was checked against the supplied snapshot on 2026-10-11. This import does not independently establish current open status or mathematical correctness. The catalogue literature assessment is dated 2026-08-22.
\subsection{Short English statement}
Identify conditions under which training homogeneous deep networks reaches a globally best normalized margin and learns beyond a fixed linearization.
\subsection{Sources}
\begin{itemize}
\item[\hypertarget{source:statement:1}{S1}] Ulam AI and the UnsolvedMath Contributors, UnsolvedMath: A Curated Collection of Open Mathematics Problems, catalogue entry. Catalogue curation and metadata: CC BY 4.0. The supplied source question is retained; definitions and the English summary are editorial additions.
\url{https://huggingface.co/datasets/ulamai/UnsolvedMath}.
\item[\hypertarget{source:statement:2}{S2}] Mini-Workshop: Mathematical Foundations of Robust and Generalizable Learning (2023), Oberwolfach Reports, DOI 10.4171/owr/2022/46, pp. 2659--2686
\url{https://doi.org/10.4171/owr/2022/46}.
\textit{Underlying Oberwolfach Reports material retains its source-specific licence, including CC BY-SA 4.0 where applicable; see THIRD\_PARTY\_NOTICES.md.}
\end{itemize}
\paragraph{Preparation and review.}
Original question attributed to the cited source and UnsolvedMath contributors. Codex (GPT-6; exact serving revision not exposed) assisted with definitions, English summaries, source inspection and formatting. No new human mathematical authorship or independent proof verification is asserted. Catalogue identity checked against the complete supplied registry and existing repository statements; known flagged repetitions and obvious identical conjectures excluded. Full mathematical-equivalence/literature review has not been performed. Statement preparation and final review followed contribution prompts C and D; automated checks verify structure and identity, not open status.
\label{end:statement}
\end{document}
